\section{Técnicas e instrumentos}
\label{sec:tecnicas_instrumentos}

Para recoger la información se usaron tres técnicas, relacionadas con la Tabla~\ref{tab:operacionalizacion}:

\begin{enumerate}
    \item \textbf{Encuesta estructurada a postulantes:} Instrumentada mediante un cuestionario en escala tipo Likert de 5 puntos (desde 1: Totalmente en desacuerdo hasta 5: Totalmente de acuerdo). El instrumento evaluó dimensiones de facilidad de navegación, claridad en los formularios de inscripción, autonomía en la recuperación de contraseñas y confianza en la reserva transparente de vacantes. La validación se realizó por juicio del tutor guía y prueba piloto con 5 estudiantes, sin requerir alfa de Cronbach por alcance descriptivo.
    \item \textbf{Entrevista semiestructurada al personal administrativo:} Conducida a partir de una guía de preguntas abiertas dirigida a secretaría y dirección institucional, con el objetivo de indagar sobre la sobrecarga laboral asociada a los reseteos manuales de contraseñas en 2025, el procedimiento de recepción y archivo de las libretas de 1ro a 3ro de secundaria, cédula de tutores, carnet de identidad y certificado de nacimiento, así como las dificultades de cruces de horarios al conformar cursos y paralelos por colegio de origen.
    \item \textbf{Pruebas técnicas y benchmarking de software:} Aplicación de pruebas de caja negra para verificar los flujos de validación de identidad en el autoservicio de contraseñas, pruebas de integridad referencial en PostgreSQL 16.4 para el registro de documentos y análisis de auditoría transaccional mediante registros de eventos (logs) de Apache 2.4 y del motor relacional. Se ejecutaron 5 corridas con Apache Benchmark 2.4.58 a $>$100 peticiones concurrentes por minuto, con tolerancia de latencia $<$1.2 s y control ACID con bloqueo pesimista \texttt{FOR UPDATE}, verificando 0 sobrecupos y estabilidad sin caídas.
\end{enumerate}
